\section{Discussion and threats to validity}
\label{sec:discussion}

\subsection{Design implications}

The contribution is an inspectable relation among proposal origin, human value authority, and successor attribution. \Cref{tab:design-checklist} turns C1--C2 into an implementation aid: map each class to persisted objects and an enforcement point, then retain a legal control and a failure probe. A check remains untested until that evidence exists. The worksheet is in \code{docs/admission-design-checklist.md}.

% Editorial design aid derived from C1--C2; no new experimental results.
\begin{table}[p]
\centering
\caption{Design checklist derived from C1--C2. Adapt the probes to the target implementation and retain persisted evidence for each outcome.}
\label{tab:design-checklist}
\begin{tabularx}{\textwidth}{@{}l>{\raggedright\arraybackslash}p{0.40\textwidth}>{\raggedright\arraybackslash}X@{}}
\toprule
Class & Retain and check & Failure probe and expected evidence \\
\midrule
$D_C$ & Exact candidate, target, evidence and producer; bind the decision at review and recheck context inside admission. & Substitute an equal-valued candidate from another context: reject without authority writes. Reload the binding to verify candidate identity. \\
\addlinespace
$D_V$ & Proposal, human-authorized value and committed value with separate roles; preserve the candidate and check the latter two values at admission. & Authorize 101 from proposal 100: commit 101, retain 100 and its identity. Reconstruction must expose false machine attribution. \\
\addlinespace
$D_F$ & Expected and observed predecessor discriminator; compare within the transaction and protect it against races. & Supersede the predecessor before replay: reject without changing the new current state. Competing admissions from one predecessor must not both commit. \\
\addlinespace
$D_B$ & Admitted field domain, actual effects and commit grouping; validate the planned successor and commit it atomically. & Change two fields; interrupt between writes. Recovery must expose all or none of that admission, never an intermediate or fragmented successor. \\
\addlinespace
$D_S$ & One resolving source per successor field; validate changed-field sources and exact unchanged-source retention before commit and on trace. & Keep values fixed but remove, duplicate or replace a source. Reject invalid construction; diagnose incomplete trace after persisted corruption. \\
\bottomrule
\end{tabularx}
\end{table}



Representations may differ while satisfying the same obligations. $D_B$ concerns the admitted change domain and complete successor; $D_S$ concerns changed-field sources and exact unchanged-source retention. Full P6 tracing also follows authorization, candidate, and evidence bindings. The practical inspection question is whether an approval names the reviewed candidate or only its value. Complete audits can reconstruct sources without specifying that review link; the checklist requires demonstrated relations, not a particular source-map representation. Its use in independent systems remains unvalidated.

\subsection{Engineering trade-offs}

Wider changes cost more on the frozen grid. The persistence-equivalent gap combines validation/planning and implementation-path differences, so it cannot price an individual check. Admission and trace require consistent canonical equality; the repair shares that relation, not a cached change set. Form-scoped tracing trades 12 database round trips for in-memory indexing and prefix traversal. The preferred access pattern remains workload-dependent.

\subsection{Formal and concrete validity}

Proposition 1 is conditional on the declared observation model. Its five class pairs, identity pair, and copy-forward control are constructions, not a general admission oracle. Alloy checks encoded scopes and selected mutations; the nine intended projections, twenty mapping mutants, catalogue, and SQLite schedules cover selected executions. The no-op and schema-uniqueness differences remain explicit in \cref{sec:conformance}.

Trace completeness checks persisted relations, not evidence truth, candidate accuracy, or human judgment. Independent oracles reduce shared-code risk but can share semantic defects, as the repaired equality bug showed. The adapter compares persisted digests and canonical locators without re-reading external evidence bytes. Detecting an out-of-band replacement that preserves metadata requires immutable/content-addressed storage or separate integrity verification.

Host authentication, faithful review presentation, serialization, hashing, policy, and adapter/transaction semantics form the trusted boundary (\cref{sec:contract}). Review-channel manipulation, indirect instructions influencing approval, and adversarial confirmation inputs were not exercised. Binding preserves a recorded decision, not its quality. Other database engines need new conformance evidence. The commit-entry policy check rejects already-visible revocation but permits later-revoked in-flight transactions (\cref{sec:authorization-timing}); it provides neither linearizable nor distributed revocation. Usability, authentication workflows, reviewer error, deletion, redaction, retention expiry, and tombstones remain untested; manual entry is outside the AI-derived claim.

\subsection{Performance and repeated-agent validity}

The empty-source comparator is a lower bound. Sequential trace timings may reflect caching, scheduling, checkpointing, and noise as well as access shape. Frozen cost results describe v8 and were not rerun after the equality repair.

The three qualified configurations from two provider families are a convenience population, not a basis for provider ranking. Endpoint names and settings do not identify immutable weights; provider load, decoding, and network conditions can vary. Pilot qualification did not ensure Final stability: D1 failures increased from 0/28 to 76/420. All planned runs remain included, with eligibility determined by \cref{tab:runtime-endpoint}.

Terminal completion and strict-trajectory compliance are distinct constructs. The post-hoc endpoint criterion admits 11 strict failures; extra calls alone establish neither caution nor general recovery competence. Original benign labels came from one author and one non-author volunteer. The first author's strict-trajectory review is a separate judgment, not an independent correction of terminal-completion labels.

The textual audit has one author coder and no human adjudication; the subsequent unblinded AI review supplies proposals only. Agreement is prevalence-sensitive and heterogeneous: A6 and B4 nearly coincide with the rule, whereas A3 contributes 32/39 disagreements. A3's unspecified target limits the construct. Author involvement, knowledge of the study purpose, and residual textual clues constrain blinding; the audit supplies human--rule comparison rather than third-party validation. Frozen rule hits remain intact.

The 720 fixed invalid-tuple calls and 179 capability-unavailable branches measure the exercised host operations. They do not sample model-generated attacks, compromised hosts or identity providers, administrator access, or model access to admission. Zero violations imply neither a population failure bound nor general prompt-injection security; zero unavailable-tool attempts reveal no behavioral heterogeneity.

\subsection{Reproducibility and evaluated-control boundary}

The deposit supports inspection without new hosted calls; reproduction in another environment is a separate test. The related-work comparison concerns disclosures through its stated cutoff, not empirical superiority. The minimal and strengthened value-audit controls both agree on seven of eight single-history cases, reconstruct field sources, and leave the paired reviewed-candidate histories indistinguishable in authority-relevant state (\cref{sec:executable-relation-result}). The evaluated safeguards therefore do not by themselves entail the exact-candidate authorization relation.

Richer review-context snapshots distinguish candidates when recorded proposal, producer, or evidence observations differ. Exact identity becomes material when distinct persisted candidates remain equivalent under those observations; it is not the only possible discriminator. This representational boundary leaves Proposition 1's conditional $D_V$ result intact: the controls test recoverability from particular audit representations, while the proposition tests omitted observation classes. Together they position an explicit correction-aware obligation, its conditional characterization, and its transactional realization. Independently developed event-sourcing/CQRS, temporal/bitemporal, and four-eyes systems were not evaluated.
